\documentclass{ximera}


\title{Adding Integers: Summary}   % Use xmlatex all -s to compile when SERVE refuses to work.
\author{Kelly Stady}
\license{CC: 0}         % replace with an appropriate license, or set it in xmPreamble

\begin{document}

\begin{abstract}

\end{abstract}

\maketitle
\label{xim:AddIntSummary}



\begin{tikzpicture}[scale=0.6]
    \begin{scope}[shift={(2,0)}]
    % Draw the owl
    \owl[transform shape,3D, body=glaucous]

    % Approximate top of head
    \coordinate (head) at (0,2.3);

    % Tilted graduation cap (using 0 rotation as base)
    \begin{scope}[rotate around={0:(head)}]
        % Cap top (diamond)
        \fill[black]
            ($(head)+(-0.9,0)$) --
            ($(head)+(0,0.3)$) --
            ($(head)+(0.9,0)$) --
            ($(head)+(0,-0.3)$) -- cycle;

        % Cap base
        \fill[black]
            ($(head)+(-0.35,-0.35)$) rectangle ($(head)+(0.35,-0.2)$);

        % --- TASSEL ---
        % 1. Small button in the center
        \fill[yellow!60!orange] (head) circle (0.05);

        % 2. Cord: Starts at center, goes to the right-most corner
        \draw[thick, orange!80!black] (head) -- ($(head)+(0.9,0)$);

        % 3. Hanging part: Drops down from that corner
        \draw[thick, orange!80!black] ($(head)+(0.9,0)$ ) -- ($(head)+(0.9,-0.6)$);

        % 4. Tassel Brush: Fanned out bottom
        \foreach \i in {-0.08, -0.04, 0, 0.04, 0.08} {
            \draw[very thin, orange!60!black] ($(head)+(0.9,-0.6)$) -- ($(head)+(0.9+\i,-0.9)$);
        }
    \end{scope}
\end{tikzpicture}


\section*{\textbf{Summary}}


\normalsize
\begin{enumerate}
\item \textbf{Adding Numbers with the Same Signs}

\begin{enumerate}
\item Add their absolute values.

\item Write the answer with the common sign.

\end{enumerate}


\item \textbf{Adding Numbers with Different Signs}

\begin{enumerate}
\item Find the absolute value of each number.

\item Subtract the smaller absolute value from the larger absolute value.

\item Write the answer using the sign of the number with the larger absolute value.

\end{enumerate}



\textbf{Don't Forget!}

These two rules came from thinking in terms of assets and debts, so you can always find an answer that way.

\end{enumerate}

\end{document}